% HAND-WRITTEN fragment -- not generated by maestro2tex, but placed by it via the
% `order` list in configs/anatop_tb_tia.json.
%
% Numbers here come from the generated tables that follow: tab_tia_amp, tab_tia_range,
% tab_tia_noise, tab_tia_drive, tab_tia_tran, tab_tia_step, tab_tia_closedloop and
% tab_tia_transfer. Re-read them from those tables if the run is repeated.
%
% Source run: Interactive.9, archived under
% simarchive/anatop_tb_tia/tia_2026-07-28/.

\subsection{Transimpedance Stage}\label{ss:tiaamp}

The transimpedance stage is amplifier A2 of the potentiostat. It holds the working
electrode (WE) at the common-mode potential $v_{\mathrm{cm}}$ on its inverting input
and converts whatever current the electrode passes into a voltage for the ADC, through
a feedback resistor $R_f$ taken from its output back to WE. Both current directions
pass through $R_f$ by Kirchhoff's current law, which is how the bipolar sensing
requirement of this revision is met without a current mirror. The transfer is

\begin{equation}\label{eq:tia-transfer}
  V_{\mathrm{OUT}} = v_{\mathrm{cm}} - I_{\mathrm{WE}} R_f ,
\end{equation}

with $I_{\mathrm{WE}}$ taken positive when current flows from the cell \emph{into} the
working-electrode node --- note that the test source in \texttt{TiaFbTB} sweeps the
opposite sign, which is why Figure~\ref{fig:tia-dc} has a positive slope. The ADC
therefore sees an offset-binary code around the $v_{\mathrm{cm}}$ code. In the
assembled channel $R_f$ is a 16-tap programmable string
(\SIrange{35}{560}{\kilo\ohm}); one instance of this stage is used per channel and all
four share the bias generator of Section~\ref{ss:biasgen-local}.

The amplifier is a two-stage Miller-compensated operational amplifier
(Figure~\ref{fig:tia-core}): a PMOS input pair on a tail current source with an NMOS
current-mirror load, then a common-source NMOS driver with a PMOS current-source load,
compensated by a MIM capacitor in series with a nulling resistor from the first-stage
output to \texttt{vout}. Two consequences of that choice run through everything below:

\begin{itemize}
\item \textbf{A PMOS input pair puts the input common-mode range against \texttt{vss}.}
The stage tracks its input to within \SI{2}{\milli\volt} from \SI{5}{\milli\volt} above
\texttt{vss} (Table~\ref{tab:tia-range}), which is what holding an electrode near ground
requires; the cost is \SI{0.44}{\volt} of unusable range at the top.
\item \textbf{The second stage is class~A, so its drive is asymmetric.} It sinks with a
driver device and sources with a fixed bias current, and the sourcing side is therefore
the limit --- \SIrange{13.6}{15.6}{\micro\ampere} against a sinking capability that this
run does not even resolve (Table~\ref{tab:tia-drive}).
\end{itemize}

The characterisation below is schematic-level, with no extracted parasitics, and comes
from four testbenches whose scope differs in ways that matter when reading the numbers:

\begin{itemize}
\item \texttt{UnityFeedbackTB} and \texttt{StepUnityFeedbackTB} close the amplifier in
unity feedback with a \SI{5}{\pico\farad} load. These give the small-signal and
large-signal figures of the \emph{amplifier}: gain, margins, offset, noise and settling.
\item \texttt{DriveTB} adds a swept DC load current to the same unity-gain
configuration, and gives the output drive limits.
\item \texttt{SlewOpenLoopTB} drives the amplifier open-loop with a rail-to-rail input
pulse, and gives the slew rates.
\item \texttt{TiaFbTB} is the only bench that closes the transimpedance loop. It wraps
$R_f = \SI{560}{\kilo\ohm}$ from output to WE and loads WE with \SI{1}{\micro\farad} of
double-layer capacitance. \textbf{Every closed-loop figure comes from here, and the
amplifier's own margins do not describe this loop} --- the electrode capacitance is
inside the feedback network.
\end{itemize}

Three limits on the coverage should be read with the results. Only the
\emph{maximum}-gain setting was simulated, and with an ideal \SI{560}{\kilo\ohm}
resistor rather than the 16-tap string, so neither the string's tap resistance nor any
lower gain setting is characterised. The run is corners only --- no Monte Carlo samples
exist yet, so offset, mirror matching and every yield figure are uncharacterised. And
the run's fifth testbench, \texttt{TypOpenLoopTB}, is deliberately not used here: its
open-loop DC operating point sits at \SI{0.23}{\volt}, with the second stage out of its
linear region, so its AC gain is not the open-loop gain. The \texttt{stb} measurement of
Figure~\ref{fig:tia-loopgain} supersedes it.
